\section*{Bab \ref{ch:b1:isostasy-gravity} --- Isostasi dan anomali gravitasi}

\begin{solution}{exo:b1:isostasy-gravity:1}
(a) $\qty{-150}{\milli\gal} = -150\times\qty{e-5}{m.s^{-2}} = \qty{-1.5e-3}{m.s^{-2}}$,
dan $\num{1.5e-3}/9.81 = \num{1.53e-4}$ dari $g$. (b) $\qty{1}{\milli\gal} =
\qty{e-5}{m.s^{-2}}$, \num{1.02e-6} dari $g$; $\qty{1}{\micro\gal} =
\qty{e-8}{m.s^{-2}}$, \num{1.02e-9} dari $g$: gravimeter yang membedakan satu mikrogal mengukur gravitasi hingga sekitar satu per miliar.
\end{solution}

\begin{solution}{exo:b1:isostasy-gravity:2}
(a) Di ekuator kedua sinus lenyap: $\gamma = \qty{9.780327}{m.s^{-2}}$. Pada
\ang{45}, $\sin^2\varphi = 0.5$ dan $\sin^2 2\varphi = 1$:
$\gamma = 9.780327\times(1 + 0.0026512 - 0.0000058) = \qty{9.806200}{m.s^{-2}}$.
Di kutub, $\sin^2\varphi = 1$ dan $\sin 2\varphi = 0$:
$\gamma = 9.780327\times1.0053024 = \qty{9.832186}{m.s^{-2}}$.
(b) $9.832186 - 9.780327 = \qty{0.051859}{m.s^{-2}}$, yaitu
\qty{5186}{\milli\gal}. (c) $3392/5186 = 0.65$: efek sentrifugal menjelaskan
sekitar dua pertiga selisih itu; kegepengan, yang menempatkan kutub lebih dekat ke
pusat, sisanya.
\end{solution}

\begin{solution}{exo:b1:isostasy-gravity:3}
$r = h\rho_c/(\rho_m - \rho_c) = 3\times2800/500 = \qty{16.8}{km}$. \omterm{def:b1:layered-earth:layers}{Kerak} menjadi
lebih tebal sebesar $h + r = \qty{19.8}{km}$ (yaitu $h\rho_m/(\rho_m - \rho_c) =
3\times3300/500$).
\end{solution}

\begin{solution}{exo:b1:isostasy-gravity:4}
(a) $1/0.3086 = \qty{3.24}{m}$. (b) $0.01/0.3086 = \qty{0.032}{m}$: ketinggian
harus diketahui hingga sekitar \qty{3}{cm}, yang menjelaskan mengapa setiap
stasiun dicatat ketinggiannya dengan penentuan posisi satelit atau sipat datar.
\end{solution}

\begin{solution}{exo:b1:isostasy-gravity:5}
(a) $\rho = 2700\times100/(100 + h)$: \qty{2673}{kg.m^{-3}} untuk \qty{1}{km},
\qty{2621}{kg.m^{-3}} untuk \qty{3}{km}, \qty{2571}{kg.m^{-3}} untuk \qty{5}{km}.
(b) $\rho = (2700\times100 - 1030\times4)/96 = \num{265880}/96 =
\qty{2770}{kg.m^{-3}}$: kolom di bawah laut lebih padat. (c) Model Pratt
menuntut perubahan lateral massa jenis pada kedalaman tetap. Di bawah dasar laut
perubahan semacam itu memang ada: \omterm{def:b1:layered-earth:lithosphere}{litosfer} muda dan panas dekat punggungan lebih ringan daripada \omterm{def:b1:layered-earth:lithosphere}{litosfer} tua dan dingin yang jauh darinya (\cref{ch:b1:heat-budget}).
Di bawah benua, sebaliknya, ketebalan \omterm{def:b1:layered-earth:layers}{kerak} terlihat berubah mengikuti relief.
\end{solution}

\begin{solution}{exo:b1:isostasy-gravity:6}
(a) Stasiun 1: $g - \gamma = 9.807751 - 9.808911 = \qty{-0.001160}{m.s^{-2}} =
\qty{-116.0}{\milli\gal}$; \omterm{def:b1:isostasy-gravity:anomalies}{koreksi udara-bebas} $0.3086\times350 =
\qty{108.0}{\milli\gal}$, sehingga $\Delta g_{\mathrm{FA}} = \qty{-8.0}{\milli\gal}$;
\omterm{def:b1:isostasy-gravity:anomalies}{koreksi Bouguer} $0.112\times350 = \qty{39.2}{\milli\gal}$, sehingga
$\Delta g_{\mathrm B} = \qty{-47.2}{\milli\gal}$.
Stasiun 2: $g - \gamma = \qty{-535.9}{\milli\gal}$; \omterm{def:b1:isostasy-gravity:anomalies}{koreksi udara-bebas}
$0.3086\times1850 = \qty{570.9}{\milli\gal}$, sehingga $\Delta g_{\mathrm{FA}} =
\qty{+35.0}{\milli\gal}$; \omterm{def:b1:isostasy-gravity:anomalies}{koreksi Bouguer} $0.112\times1850 =
\qty{207.2}{\milli\gal}$, sehingga $\Delta g_{\mathrm B} = \qty{-172.2}{\milli\gal}$.
(b) Di stasiun 2 \omterm{def:b1:isostasy-gravity:anomalies}{anomali Bouguer} sangat negatif: di bawah pegunungan ada defisit
massa, akar dari \omterm{def:b1:layered-earth:layers}{kerak} ringan. \omterm{def:b1:isostasy-gravity:anomalies}{Anomali udara-bebas}, kecil dibanding \omterm{def:b1:isostasy-gravity:anomalies}{koreksi Bouguer},
menunjukkan bahwa defisit itu hampir menyeimbangkan pegunungan: pegunungan itu sebagian besar terkompensasi
(sampai sekitar $172/207$, atau \qty{83}{\percent}, bila model lempeng berlaku untuk seluruh
rentang).
\end{solution}

\begin{solution}{exo:b1:isostasy-gravity:7}
(a) Seperti dalam pembuktian proposisi, lapisan tipis pada jarak $z$
menyumbang $G\rho\,2\pi z\,h\int_0^a s\,(s^2 + z^2)^{-3/2}\,\mathrm ds$, kini
dengan integral dihentikan pada $s = a$:
$\int_0^a s\,(s^2 + z^2)^{-3/2}\,\mathrm ds = 1/z - 1/\sqrt{z^2 + a^2}$. Maka
$g = 2\pi G\rho h\,\bigl(1 - z/\sqrt{z^2 + a^2}\bigr)$, yang menuju nilai lempeng
bila $a \to \infty$. (b) $z/\sqrt{z^2 + a^2} = 0.1$ memberikan
$\sqrt{z^2 + a^2} = 10z$, $a = \sqrt{99}\,z = \qty{9.95}{km}$ untuk
$z = \qty{1}{km}$. (c) Batuan dalam kira-kira sepuluh kali kedalamannya di bawah
stasiun memberikan \qty{90}{\percent} dari tarikan itu: koreksi medan
memerlukan relief dalam radius beberapa kilometer hingga puluhan kilometer di sekitar setiap stasiun, bukan hanya di bawahnya.
\end{solution}

\begin{solution}{exo:b1:isostasy-gravity:8}
(a) $r' = d(\rho_c - \rho_w)/(\rho_m - \rho_c) = 4\times1770/500 =
\qty{14.16}{km}$, dan \omterm{def:b1:layered-earth:layers}{kerak} menjadi $35 - 4 - 14.16 = \qty{16.84}{km}$ tebalnya.
(b) Model Airy dengan satu massa jenis \omterm{def:b1:layered-earth:layers}{kerak} memprediksi \omterm{def:b1:layered-earth:layers}{kerak} samudra yang lebih
tebal lebih dari tiga kali lipat. \omterm{def:b1:layered-earth:layers}{Kerak} samudra bukanlah \omterm{def:b1:layered-earth:layers}{kerak} benua yang menipis: ia adalah batuan lain yang lebih padat, dan kedua tingkat Bumi adalah dua jenis \omterm{def:b1:layered-earth:layers}{kerak},
seperti disarankan oleh \omterm{def:b1:isostasy-gravity:hypsometry}{kurva hipsometrik}.
\end{solution}

\begin{solution}{exo:b1:isostasy-gravity:9}
(a) Satu detik busur adalah $\pi/(180\times3600) = 1/\num{206265}$ radian:
\qty{5.236}{''} $= \qty{2.538e-5}{rad}$ dan \qty{15.885}{''} $=
\qty{7.701e-5}{rad}$. (b) $a_h = g\theta$: $9.79\times\num{2.538e-5} =
\qty{2.485e-4}{m.s^{-2}} = \qty{24.9}{\milli\gal}$ teramati, dan
$9.79\times\num{7.701e-5} = \qty{75.4}{\milli\gal}$ dihitung dari gunung yang terlihat.
(c) $1 - 5.236/15.885 = 0.670$: kompensasi membatalkan sekitar dua pertiga
tarikan gunung yang terlihat.
\end{solution}

\begin{solution}{exo:b1:isostasy-gravity:10}
(a) Di luarnya, bola menarik seperti massa titik di pusatnya, dengan
$GM/(x^2 + z^2)$ sepanjang garis ke pusat; komponen vertikalnya adalah itu dikali kosinus
$z/\sqrt{x^2 + z^2}$, yang memberikan $GMz/(x^2 + z^2)^{3/2}$.
(b) $M = \tfrac43\pi(\qty{2000}{m})^3\times(\qty{-200}{kg.m^{-3}}) =
\qty{-6.70e12}{kg}$, dan di atas pusat $GM/z^2 = \num{6.674e-11}\times
(\num{-6.70e12})/(4000)^2 = \qty{-2.80e-5}{m.s^{-2}} = \qty{-2.80}{\milli\gal}$.
(c) $(1 + x^2/z^2)^{-3/2} = 1/2$ memberikan $x = z\sqrt{2^{2/3} - 1} = 0.766\,z =
\qty{3.07}{km}$. (d) Lebar setengah-maksimum profil memberikan kedalaman:
$z = 1.305\,x_{1/2}$. Benda yang lebih dalam memberikan anomali yang lebih lebar dan lebih datar.
\end{solution}

\begin{solution}{exo:b1:isostasy-gravity:11}
(a) $V = \tfrac13\pi R^2 H = \tfrac13\pi(\qty{15e3}{m})^2(\qty{3e3}{m}) =
\qty{7.07e11}{m^3}$, dan $\Delta m = (2800 - 1030)V = \qty{1.25e15}{kg}$.
(b) Pusat massa berada \qty{0.75}{km} di atas dasar, pada
$d = 5 - 0.75 = \qty{4.25}{km}$ di bawah permukaan laut. Di sana anomali potensial
adalah $T = -G\Delta m/d$, sehingga $N = G\Delta m/(\gamma d) =
\num{6.674e-11}\times\num{1.25e15}/(9.81\times4250) = \qty{2.0}{m}$.
(c) Sebagian besar massa kerucut berada di bagian bawahnya yang lebar, pada 5 hingga
\qty{16}{km} dari titik di atas puncak: rata-rata $1/r$ atas massa lebih kecil daripada
$1/\qty{4.25}{km}$, dan massa titik melebih-lebihkan potensialnya. Menjumlahkan
cakram tipis memberikan sekitar \qty{1.1}{m}.
(d) Akar \omterm{def:b1:layered-earth:layers}{kerak} ringan di bawahnya akan menempatkan defisit dengan massa yang sama
lebih dalam; kedua potensial hampir saling meniadakan, dan tonjolannya akan jauh lebih kecil,
tetapi tidak nol, karena defisit terletak lebih jauh dari permukaan daripada kelebihan massa.
\end{solution}

\begin{solution}{exo:b1:isostasy-gravity:12}
(a) Laju pengangkatannya adalah $-\mathrm dw/\mathrm dt = w_0\,e^{-t/\tau}/\tau = w/\tau$.
(b) $w = 818\,e^{-10000/4000} = 818\times0.0821 = \qty{67.1}{m}$, dan laju
$\qty{67.1}{m}/\num{4000}\,\unit{yr} = \qty{16.8}{mm.yr^{-1}}$, lebih dari satu
setengah kali nilai terukur. (c) Dengan $\tau = 3100$ tahun, $w =
818\,e^{-3.226} = \qty{32.5}{m}$ dan lajunya \qty{10.5}{mm.yr^{-1}}; dengan
$\tau = \num{68000}$ tahun, $w = 818\,e^{-0.147} = \qty{706}{m}$ dan lajunya
\qty{10.4}{mm.yr^{-1}}. Keduanya menyiratkan depresi tersisa yang sangat berbeda.
Depresi yang tidak dikompensasi adalah defisit massa batuan \omterm{def:b1:layered-earth:layers}{mantel}: pada wilayah
luas ia memberikan \omterm{def:b1:isostasy-gravity:anomalies}{anomali udara-bebas} sekitar $-2\pi G\rho_m w$, yaitu
$-0.0419\times3.3\times706 \approx \qty{-98}{\milli\gal}$ untuk kasus lambat,
dibanding \qty{-4}{\milli\gal} untuk yang cepat. Survei gravitasi (atau ketinggian
garis pantai lama) membedakan keduanya. (d) Konstanta waktu mengukur seberapa cepat
\omterm{def:b1:layered-earth:layers}{mantel} mengalir kembali di bawah depresi: viskositasnya.
\end{solution}

\begin{solution}{pb:b1:isostasy-gravity:1}
\textbf{1.} $r = h\rho_c/(\rho_m - \rho_c) = 4\times2700/600 = \qty{18}{km}$.

\textbf{2.} \omterm{def:b1:layered-earth:layers}{Kerak} setebal $35 + 4 + 18 = \qty{57}{km}$; dasarnya, \omterm{def:b1:layered-earth:discontinuity}{Moho},
berada $35 + 18 = \qty{53}{km}$ di bawah permukaan laut.

\textbf{3.} Dataran tinggi: \qty{57}{km} \omterm{def:b1:layered-earth:layers}{kerak} dan $60 - 53 = \qty{7}{km}$ \omterm{def:b1:layered-earth:layers}{mantel},
$\num{57000}\times2700 + \num{7000}\times3300 = \num{1.539e8} + \num{2.31e7} =
\qty{1.770e8}{kg.m^{-2}}$. Dataran rendah: $\num{35000}\times2700 + \num{25000}\times3300
= \num{9.45e7} + \num{8.25e7} = \qty{1.770e8}{kg.m^{-2}}$. Keduanya sama, seperti
\cref{prop:b1:isostasy-gravity:equal-mass} mensyaratkannya.

\textbf{4.} $p = g\times\text{massa per satuan luas} = 9.81\times\num{1.770e8} =
\qty{1.74e9}{Pa}$, sekitar \qty{1.7}{GPa}, di bawah kedua kolom.

\textbf{5.} Akar harus mengimbangi, dengan massa jenisnya yang lebih kecil, massa dataran tinggi:
setiap meter akar menghilangkan $\rho_m - \rho_c = \qty{600}{kg.m^{-3}}$ massa,
dan setiap meter dataran tinggi menambahkan \qty{2700}{kg.m^{-3}}, sehingga $2700/600 = 4.5$
meter akar diperlukan. Dengan \omterm{def:b1:layered-earth:layers}{kerak} \qty{2900}{kg.m^{-3}}, rasionya
akan menjadi $2900/400 = 7.25$: semakin kecil kontrasnya, semakin dalam akarnya.

\textbf{6.} \omterm{def:b1:layered-earth:discontinuity}{Moho}, lompatan kecepatan seismik yang kuat, akan ditemukan
semakin dalam di bawah tepi dataran tinggi, dari \qty{35}{km} di bawah dataran rendah
hingga sekitar \qty{53}{km} di bawah dataran tinggi.

\textbf{7.} $\rho = 2700\times100/104 = \qty{2596}{kg.m^{-3}}$.

\textbf{8.} $\num{104000}\times2596.2 = \qty{2.700e8}{kg.m^{-2}}$ dan
$\num{100000}\times2700 = \qty{2.700e8}{kg.m^{-2}}$.

\textbf{9.} $2700 - 2596 = \qty{104}{kg.m^{-3}}$, atau $104/2700 =
\qty{3.85}{\percent}$.

\textbf{10.} $\Delta\rho/\rho = \alpha\,\Delta T$ memberikan $\Delta T =
0.0385/\num{3e-5} = \qty{1280}{K}$ rata-rata di seluruh kolom \qty{104}{km}. Itu akan melelehkan sebagian besar kolom: panas saja tidak dapat membuat dataran tinggi ini
cukup ringan, dan mekanisme Pratt tidak dapat menjadi yang utama di sini.

\textbf{11.} Kolom Pratt tidak memiliki akar: \omterm{def:b1:layered-earth:discontinuity}{Moho} tetap sekitar \qty{35}{km}
di bawah dataran tinggi, seperti di bawah dataran rendah.

\textbf{12.} \qty{52}{km} dekat dengan \qty{53}{km} milik Airy dan jauh dari
\qty{35}{km} milik Pratt: dataran tinggi terkompensasi terutama oleh akar \omterm{def:b1:layered-earth:layers}{kerak}.

\textbf{13.} $0.0419\times2.7\times4000 = \qty{452.5}{\milli\gal}$.

\textbf{14.} Tidak terkompensasi (\cref{prop:b1:isostasy-gravity:broad}):
$\Delta g_{\mathrm{FA}} \approx \qty{+452.5}{\milli\gal}$,
$\Delta g_{\mathrm B} \approx 0$.

\textbf{15.} Terkompensasi: $\Delta g_{\mathrm{FA}} \approx 0$ dan
$\Delta g_{\mathrm B} \approx \qty{-452.5}{\milli\gal}$. Massa akar yang hilang
memberikan $-0.0419\times0.6\times\num{18000} = \qty{-452.5}{\milli\gal}$, sama,
karena $\Delta\rho\,r = \rho_c h$.

\textbf{16.} $\sin^2\ang{30} = 0.25$ dan $\sin^2\ang{60} = 0.75$:
$\gamma = 9.780327\times(1 + 0.0013256 - 0.0000044) = \qty{9.793249}{m.s^{-2}}$.
Lalu $g - \gamma = 9.781085 - 9.793249 = \qty{-1216.4}{\milli\gal}$; \omterm{def:b1:isostasy-gravity:anomalies}{koreksi udara-bebas}
adalah $0.3086\times4000 = \qty{1234.4}{\milli\gal}$, sehingga
$\Delta g_{\mathrm{FA}} = \qty{+18.0}{\milli\gal}$, dan
$\Delta g_{\mathrm B} = 18.0 - 452.5 = \qty{-434.5}{\milli\gal}$.

\textbf{17.} $434.5/452.5 = 0.96$: dataran tinggi sekitar \qty{96}{\percent} terkompensasi
(\omterm{def:b1:isostasy-gravity:anomalies}{anomali udara-bebas} positif yang kecil adalah bagian yang tidak terkompensasi,
\qty{4}{\percent}).

\textbf{18.} Tidak. Dekat pusat dataran tinggi yang lebar setiap lapisan anomali bertindak
sebagai lempeng, yang tarikannya tidak bergantung pada kedalamannya: hanya defisit massa
per satuan luas yang berpengaruh, dan itu sama, $\rho_c h$, di kedua model. Keduanya
berbeda hanya di dekat tepi, tempat kedalaman massa kompensasi mengubah
bentuk anomali.

\textbf{19.} Jauh dari tepi di kedua sisi, \omterm{def:b1:isostasy-gravity:anomalies}{anomali udara-bebas} mendekati nol.
Beberapa kilometer di dalam tepi, dataran tinggi di bawah stasiun dekat dan
meluas jauh hanya di satu sisi: ia tetap menarik dengan hampir seluruh nilai lempengnya.
Akar, \qty{35}{km} atau lebih di bawah, dari situ tampak sebagai benda yang berakhir
hampir lurus di bawah: ia menarik dengan sedikit lebih dari separuh pengaruhnya. Selisihnya
adalah puncak positif. Beberapa kilometer di luar, dataran tinggi menarik
sangat sedikit, sementara akar tetap menarik dengan hampir separuh pengaruhnya: palung negatif.
Profilnya adalah pasangan positif--negatif melintasi setiap tepi, seperti dalam
gambar bab ini.

\textbf{20.} \omterm{def:b1:layered-earth:layers}{Kerak} kini $57 - 1 = \qty{56}{km}$ tebalnya, $\qty{21}{km}$ lebih
daripada dataran rendah. Dalam kesetimbangan kelebihan ini terbagi menjadi $h' + r'$ dengan
$r' = 4.5\,h'$: $h' = 21/5.5 = \qty{3.818}{km}$ dan $r' = \qty{17.18}{km}$.

\textbf{21.} Dataran tinggi kini berdiri pada \qty{3818}{m}.

\textbf{22.} Batuan yang kini di permukaan berada \qty{1}{km} di bawah permukaan lama,
pada \qty{3000}{m}; kini berada pada \qty{3818}{m}: ia telah naik sebesar \qty{818}{m},
yaitu $1000\times\rho_c/\rho_m$.

\textbf{23.} Setiap kilometer yang dierosi menurunkan permukaan sebesar $1 - 2700/3300 =
0.182$ kilometer, sehingga menurunkannya \qty{4}{km} memerlukan
$4/0.182 = \qty{22}{km}$ erosi, seluruh kelebihan ketebalan $h + r$. Pada
sabuk gunung tua yang telah aus, batuan di permukaan karena itu pernah berada kira-kira
dua puluh kilometer dalamnya: akar dalam sabuk itu tersingkap.

\textbf{24.} Permukaan turun sebesar $1000 - 818$: untuk setiap kilometer batuan
yang dierosi, penurunan bersih permukaan adalah
\textbf{\qty[text-series-to-math=true]{182}{m}}.
\end{solution}
